\documentclass[10pt,letterpaper]{article}
\usepackage{technical-testing}
\setDocTitle{Chapter-to-Competency Crosswalk}
\setDocSubtitle{Traceability between book concepts and engineering evidence}
\setDocVersion{1.0}
\setDocStatus{Educational / Draft}
\setDocOwner{Game Development}
\setDocDate{2026-10-30}
\setDocSummary{Traceability between book concepts and engineering evidence}
\begin{document}
\makedoctitle
\makedocsummary
\tableofcontents\clearpage
\section{Competency model}
There are four independent evidence streams: \textbf{reading} (chapter concepts), \textbf{visual design} (poses, staging, observation), \textbf{engineering} (implementation and tests), and \textbf{evaluation} (playtests and measurement). Completion in one stream does not imply mastery of the others.
\section{Module to evidence mapping}
\begin{longtable}{p{.27\linewidth}p{.19\linewidth}p{.38\linewidth}}
\textbf{Module}&\textbf{Book chapters}&\textbf{Observable evidence}\\\hline
Direction and timing communication&1--17&Storyboard, timing sheet, reproducible timing handoff\\
Force and physical motion&18--26&Arc predictions, force/weight comparative shots\\
Expressive spacing and cycles&27--46&Frame strips, holds, overlap analysis, seamless loops\\
Teams and procedural effects&47--54&Deterministic varied agent motion and effects timing\\
Locomotion and acting&55--67&Walk/run contact tracks, characterization exercises\\
Dialogue and music&68--72&Viseme annotation, sound-marker synchronization\\
Playable animation and editorial&73--82&State graph, runtime data contract, edited playable sequence\\
\end{longtable}
\section{Crosswalk rules}
One book chapter may inform several labs. Mark as \emph{Source-aligned} only claims or exercises rooted in the chapter's actual subject. Label later topics (root motion reconciliation, deterministic accumulator loops, engine graphs, CPU budgets, tests) as \emph{Supplementary engineering}. Never imply the original publication specifies a named modern engine API.
\section{Pass criteria}
Each evidence item states: objective, technical constraints, measurement method, captured output, source chapter, evaluator, date, deviation and next action. Score observed intent and implementation correctness separately.
\end{document}
